% =====================================================================
%  1.2 Kvadratsetningene og faktorisering — fasit
%  \fk{overskrift}   \fo{oppgavenummer} svar
% =====================================================================

\fk{Del 1 — Kvadratsetningene}
\fo{1.1} a) $x^2 + 10x + 25$ \ b) $a^2 - 6a + 9$ \ c) $16b^2 + 8b + 1$ \ d) $m^2 + 2mn + n^2$
\fo{1.2} a) $x^2 - 16$ \ b) $a^2 - 49$ \ c) $4p^2 - 9$ \ d) $25 - y^2$
\fo{1.3} a) $2ab$ \ b) $b^2$
\fo{1.4} a) $12x$ \ b) $4a^2 + 1$ \ c) $4x^2 - 4$ \ d) $4m + 8$
\fo{1.5} a) $2601$ \ b) $9604$ \ c) $899$
\fo{1.6} Feilen er i overgangen til $-(x^2 - 8x - 16)$: $(x-4)^2 = x^2 - 8x + 16$. Riktig: $x^2 + 2x + 1 - x^2 + 8x - 16 = 10x - 15$.
\fo{1.7} a) Alltid. $(b-a)^2 = (a-b)^2$, fordi $(-1)^2 = 1$.
b) Noen ganger. Forskjellen er $4x$, så likheten gjelder bare når $x = 0$.
c) Aldri. $(x-5)^2 = x^2 - 10x + 25$, så forskjellen er alltid $50$.
\fo{1.8} $(x+a)^2 = x^2 + 2ax + a^2$. Da er $2a = 10$, så $a = 5$ og $b = a^2 = 25$.
\fo{1.9} a) $(a+b)^2 - (a-b)^2 = (a^2 + 2ab + b^2) - (a^2 - 2ab + b^2) = 4ab$.
b) $103^2 - 97^2 = 4 \cdot 100 \cdot 3 = 1200$.
\fo{Sjekk 1} a) $x^2 - 4x + 4$ \ b) $21$ \ c) Et talleksempel: $a = 1$ og $b = 1$ gir $(1+1)^2 = 4$, men $1^2 + 1^2 = 2$. Arealfiguren viser at $2ab$ mangler.

\fk{Del 2 — Faktorisering}
\fo{2.1} a) $3(x+3)$ \ b) $4(a-2b)$ \ c) $x(x+1)$ \ d) $3m(2m-1)$
\fo{2.2} a) $(x+2)(x-2)$ \ b) $(a+5)(a-5)$ \ c) $(3+y)(3-y)$
\fo{2.3} a) $(x+4)^2$ \ b) $(2x+3)^2$
\fo{2.4} a) $3(x^2 - 4x + 4) = 3(x-2)^2$ \ b) $2a(a-2)(a+2)$ \ c) $3(x-2)(x+2)$ \ d) $x(x-1)(x+1)$
\fo{2.5} a) Nei. $x$ kan trekkes ut av $x^2 + 2x$, så faktoriseringen er ikke fullstendig.
b) $3a^2(2a-3)$
\fo{2.6} $x - 3$ cm. Kontroll: $(x-3)(x+3) = x^2 - 9$.
\fo{2.7} a) Noen ganger. $(x+4)^2 = x^2 + 8x + 16$, så likheten gir $8x = 0$, altså bare $x = 0$.
b) Aldri. $(x+2)(x-2) = x^2 - 4$, og $x^2 - 4 \neq x^2 + 4$ for alle $x$.
c) Alltid. $(3x-2)(3x+2) = (3x)^2 - 2^2 = 9x^2 - 4$.
\fo{2.8} Mange svar er riktige. Eksempel: $2x^2 - 50 = 2(x^2 - 25) = 2(x-5)(x+5)$. Annet eksempel: $2(x+1)^2 - 72 = 2(x-5)(x+7)$ ($(x-5)$ er faktor uten $(x+5)$). Forklaring: $2x^2 + 50 = 2(x^2 + 25)$, og $x^2 + 25$ er en sum av kvadrater, så konjugatsetningen kan ikke brukes. Dessuten er $(x-5)$ ikke en faktor, for $x = 5$ gir $100 \neq 0$.
\fo{Sjekk 2} a) $3a(a-2)$ \ b) $(x+7)(x-7)$ \ c) $2(x+2)^2$

\fk{Ekstra utfordring}
\fo{U1} $(x+y+1)^2 = x^2 + y^2 + 1 + 2xy + 2x + 2y$.
\fo{U2} a) $3x(x-3)(x+3)$ \ b) $3(2a-b)(2a+b)$
\fo{U3} $(n+1)^2 - n^2 = n^2 + 2n + 1 - n^2 = 2n + 1$. Og $2n + 1$ er et oddetall for alle heltall $n$, fordi $2n$ er et partall.
\fo{U4} Fra $(x+c)^2 = x^2 + 2cx + c^2$ må $c^2 = 25$, så $c = 5$ eller $c = -5$ og $k = 2c = \pm 10$. Da er $x^2 + 10x + 25 = (x+5)^2$ og $x^2 - 10x + 25 = (x-5)^2$.
For $x^2 + 6x + k$ må $2c = 6$, så $c = 3$ og $k = 9$. Første gang bestemmer $c^2 = 25$ to verdier av $c$. Andre gang bestemmer $2c = 6$ bare én.
\fo{U5} a) Programmet skriver ut \texttt{1 True} til \texttt{10 True}. Alle linjene er sanne.
b) Nei. Programmet sjekker bare $n = 1, \ldots, 10$. Et bevis må gjelde for alle heltall $n$, og det gjør den algebraiske utregningen i U3.
\fo{U6} a) $(m-d)(m+d) = m^2 + md - dm - d^2 = m^2 - d^2$.
b) $47 \cdot 53 = (50-3)(50+3) = 50^2 - 3^2 = 2500 - 9 = 2491$.
c) Tallene er $m - d$ og $m + d$, med sum $2m = 100$, så $m = 50$. Produktet er $50^2 - d^2 = 2500 - d^2$. Det er størst når $d = 0$. Svar: $50$ og $50$, produkt $2500$.
\fo{U7} a) $x^4 - 16 = (x^2-4)(x^2+4) = (x-2)(x+2)(x^2+4)$.
b) Hvis $x^2 + 4 = (x-r)(\ldots)$, ville uttrykket vært $0$ for $x = r$. Men $x^2 + 4 \geq 4$ for alle reelle $x$.
\fo{U8} Disse går: $1, 3, 4, 5, 7, 8, 9, 11, 12, 13, 15, 16, 17, 19, 20$. Disse går ikke: $2, 6, 10, 14, 18$. De har rest $2$ ved deling med $4$.
Begrunnelse: $a^2 - b^2 = (a-b)(a+b)$. Faktorene $a-b$ og $a+b$ har differanse $2b$, så de er begge oddetall eller begge partall. Begge odde gir et oddetall. Begge partall gir et tall delelig med $4$. Rest $2$ er umulig.
